\documentclass[11pt]{article}
\usepackage{amsmath,amssymb,amsthm}
\usepackage[margin=2.5cm]{geometry}

\newtheorem{theorem}{Theorem}
\newtheorem{lemma}[theorem]{Lemma}
\theoremstyle{definition}
\newtheorem{definition}[theorem]{Definition}
\theoremstyle{remark}
\newtheorem{remark}[theorem]{Remark}
\theoremstyle{definition}
\newtheorem*{statement}{Official statement}

\newcommand{\T}{\mathbb{T}}

\title{Problem 1, Cell 1: the planar case $d=2$}
\author{}
\date{}

\begin{document}
\maketitle

\begin{statement}
A line here always means a line through the origin of $\mathbb{R}^d$. The angle between two lines
$\ell,\ell'$ is the acute (non-obtuse) angle $\theta(\ell,\ell')\in[0,\pi/2]$ between them. If
$\ell,\ell'$ are spanned by unit vectors $x,x'$, then $\theta(\ell,\ell')=\arccos|\langle x,x'\rangle|$.
For lines $\ell_1,\dots,\ell_N$ in $\mathbb{R}^d$ (repetitions allowed), write
$S(\ell_1,\dots,\ell_N)=\sum_{1\le i<j\le N}\theta(\ell_i,\ell_j)$.

We begin in the plane ($d=2$), where the conjectured optimum splits the $N$ lines as evenly as
possible between two perpendicular directions. Let $\ell_1,\dots,\ell_N$ be lines in $\mathbb{R}^2$.
Prove that $S(\ell_1,\dots,\ell_N)\le(\pi/2)\cdot\lfloor N^2/4\rfloor$.
\end{statement}

\noindent Theorem~\ref{thm:main} below is exactly this statement. It is proved for arbitrary lines in
$\mathbb{R}^2$, and lines through the origin are a special case. For lines spanned by unit vectors
$x,x'$, the angle $\angle$ used below equals $\arccos|\langle x,x'\rangle|=\theta$. Repeated lines are
allowed.

\paragraph{Known result.} The planar case is known~\cite{FVZ16,BM19}; the proof below is
self-contained.

\paragraph{Setting.}
Let $\ell_1,\dots,\ell_N$ be lines in $\mathbb{R}^2$. The angle $\angle(\ell_i,\ell_j)\in[0,\pi/2]$
between two lines is the smaller of the two angles formed by their directions (it is $0$ for
parallel lines, and it depends only on the directions of the lines). Put
\[
  S=\sum_{1\le i<j\le N}\angle(\ell_i,\ell_j).
\]

\begin{theorem}\label{thm:main}
$S\le \dfrac{\pi}{2}\left\lfloor \dfrac{N^2}{4}\right\rfloor$.
\end{theorem}

\section{Directions as points of a circle of length $\pi$}

Let $\T=\mathbb{R}/\pi\mathbb{Z}$, with the Lebesgue measure $\lambda$ inherited from $[0,\pi)$
(so $\lambda(\T)=\pi$) and the quotient map $q\colon\mathbb{R}\to\T$. The direction of a line
is determined by an angle $\varphi\in\mathbb{R}$ with the $x$-axis, up to adding multiples of
$\pi$; hence each $\ell_i$ has a well-defined direction $\varphi_i\in\T$.

For $u\in\T$ define $\|u\|=\min\{|t| : t\in\mathbb{R},\ q(t)=u\}\in[0,\pi/2]$. If $t\in[0,\pi)$ is
the representative of $u$, then $\|u\|=\min(t,\pi-t)$. By definition of the angle between two
lines,
\begin{equation}\label{eq:angle}
  \angle(\ell_i,\ell_j)=\|\varphi_i-\varphi_j\|.
\end{equation}

\begin{remark}\label{rem:fundamental}
Any half-open interval $[s,s+\pi)\subset\mathbb{R}$ is mapped by $q$ bijectively onto $\T$, and
$q$ restricted to it preserves the measure of Borel sets. For $s=0$ this is the definition of
$\lambda$, and for general $s$ it follows by translation invariance of the Lebesgue measure. This is
the only fact about
``wrap-around'' we shall need.
\end{remark}

\section{The arcs}

For each $i$ choose a real representative $p_i$ of $\varphi_i$ and set
\[
  I_i = q\bigl([\,p_i-\tfrac{\pi}{4},\ p_i+\tfrac{\pi}{4})\bigr)\subset\T .
\]
This set does not depend on the choice of $p_i$ (changing $p_i$ by $k\pi$ translates the interval
by $k\pi$). By Remark~\ref{rem:fundamental}, $\lambda(I_i)=\pi/2$, since
$[p_i-\frac\pi4,p_i+\frac\pi4)$ is contained in a half-open interval of length $\pi$ on which $q$
is injective and measure preserving. (This value is not used below; it only shows that each arc
covers exactly half of $\T$.) We write $\mathbf 1_i$ for the indicator function of $I_i$.

\begin{lemma}[Symmetric difference]\label{lem:symdiff}
For all $i\neq j$,
\[
  \lambda(I_i\,\triangle\, I_j)=\int_{\T}\bigl|\mathbf 1_i(x)-\mathbf 1_j(x)\bigr|\,dx
  = 2\,\angle(\ell_i,\ell_j).
\]
\end{lemma}

\begin{proof}
The first equality holds because $|\mathbf 1_i-\mathbf 1_j|$ is the indicator function of
$I_i\triangle I_j$. For the second, let $\theta=\angle(\ell_i,\ell_j)=\|\varphi_j-\varphi_i\|\in[0,\pi/2]$.
Let $t\in[0,\pi)$ be the representative of $\varphi_j-\varphi_i$. Then $\theta=t$ or $\theta=\pi-t$.
In the second case, $\varphi_i-\varphi_j$ has representative $\pi-t=\theta$ (if $t>0$; if $t=0$ the
first case holds). Since the statement is symmetric in $i$ and $j$, we may swap the two indices
if needed and assume
\[
  \varphi_j-\varphi_i=q(\theta),\qquad 0\le\theta\le\pi/2 .
\]
Choose the representatives $p_i$ and $p_j=p_i+\theta$. Let $J_i=[p_i-\frac\pi4,p_i+\frac\pi4)$ and
$J_j=[p_i+\theta-\frac\pi4,p_i+\theta+\frac\pi4)$, so that $I_i=q(J_i)$ and $I_j=q(J_j)$.

\emph{No wrap-around.} Because $0\le\theta\le\pi/2$, both $J_i$ and $J_j$ lie in the half-open
interval $F=[p_i-\frac\pi4,\ p_i+\frac{3\pi}4)$, which has length exactly $\pi$. (Here we use
$\theta+\frac\pi4\le\frac{3\pi}4$.) By Remark~\ref{rem:fundamental}, $q|_F$ is a measure-preserving
bijection $F\to\T$. Being a bijection, it commutes with set operations:
$q(J_i)\triangle q(J_j)=q(J_i\triangle J_j)$. Hence
$\lambda(I_i\triangle I_j)$ equals the Lebesgue measure of $J_i\triangle J_j$ computed in $\mathbb{R}$.

\emph{Computation in $\mathbb{R}$.} Since $\theta\le\pi/2$, the left end of $J_j$ satisfies
$p_i+\theta-\frac\pi4\le p_i+\frac\pi4$, so
\[
  J_i\setminus J_j=[\,p_i-\tfrac\pi4,\ p_i+\theta-\tfrac\pi4),\qquad
  J_j\setminus J_i=[\,p_i+\tfrac\pi4,\ p_i+\theta+\tfrac\pi4).
\]
Each has length $\theta$. (If $\theta=\pi/2$ the two intervals are disjoint and adjacent, and the
formulas still hold. If $\theta=0$ both sets are empty.) Therefore
$\lambda(I_i\triangle I_j)=2\theta$.
\end{proof}

\section{Counting pairs}

Define the covering function
\[
  A(x)=\#\{\,i\in\{1,\dots,N\} : x\in I_i\,\}=\sum_{i=1}^N\mathbf 1_i(x),\qquad x\in\T .
\]
It is a finite sum of indicator functions of arcs, so it is a bounded measurable (step) function
with integer values in $\{0,1,\dots,N\}$.

\begin{lemma}[Pair count]\label{lem:count}
For every $x\in\T$,
\[
  \sum_{1\le i<j\le N}\bigl|\mathbf 1_i(x)-\mathbf 1_j(x)\bigr| = A(x)\,\bigl(N-A(x)\bigr).
\]
\end{lemma}

\begin{proof}
Fix $x$. The number $|\mathbf 1_i(x)-\mathbf 1_j(x)|$ is $1$ if exactly one of $I_i,I_j$ contains
$x$, and $0$ otherwise. So the left-hand side counts unordered pairs $\{i,j\}$ with exactly one
of the arcs containing $x$. Let $P=\{i : x\in I_i\}$, so $|P|=A(x)$. Such a pair consists of one
index in $P$ and one index outside $P$. Each unordered pair of this kind corresponds to exactly
one element of $P\times(\{1,\dots,N\}\setminus P)$, and conversely. Hence there are
$A(x)(N-A(x))$ such pairs.
\end{proof}

\begin{lemma}\label{lem:int}
For every integer $a$ with $0\le a\le N$ we have $a(N-a)\le\lfloor N^2/4\rfloor$.
\end{lemma}

\begin{proof}
$a(N-a)=\frac{N^2}{4}-\bigl(a-\frac N2\bigr)^2$. If $N$ is even this is at most $N^2/4=\lfloor N^2/4\rfloor$.
If $N$ is odd, $a-\frac N2$ is a half-integer, so $(a-\frac N2)^2\ge\frac14$ and
$a(N-a)\le\frac{N^2-1}{4}=\lfloor N^2/4\rfloor$.
\end{proof}

\section{Proof of Theorem~\ref{thm:main}}

By \eqref{eq:angle}, Lemma~\ref{lem:symdiff}, linearity of the integral (a finite sum), and
Lemma~\ref{lem:count},
\[
  2S=\sum_{i<j}2\angle(\ell_i,\ell_j)
    =\sum_{i<j}\int_{\T}|\mathbf 1_i-\mathbf 1_j|\,dx
    =\int_{\T}\sum_{i<j}|\mathbf 1_i-\mathbf 1_j|\,dx
    =\int_{\T}A(x)\bigl(N-A(x)\bigr)\,dx .
\]
By Lemma~\ref{lem:int} the integrand is at most $\lfloor N^2/4\rfloor$ at every point, and
$\lambda(\T)=\pi$. Hence $2S\le\pi\lfloor N^2/4\rfloor$, that is,
$S\le\frac\pi2\lfloor N^2/4\rfloor$. \qed

\begin{remark}[Sharpness]
Take $\lfloor N/2\rfloor$ lines parallel to the $x$-axis and $\lceil N/2\rceil$ lines parallel to the
$y$-axis. Exactly $\lfloor N/2\rfloor\lceil N/2\rceil=\lfloor N^2/4\rfloor$ pairs are
perpendicular, and all other pairs have angle $0$. So equality holds and the bound is best possible.
\end{remark}

\begin{thebibliography}{9}
\bibitem{BM19} D.~Bilyk and R.~W.~Matzke, \emph{On the Fejes T\'oth problem about the sum of angles
  between lines}, Proc.\ Amer.\ Math.\ Soc.\ \textbf{147} (2019); arXiv:1801.07837.
\bibitem{FVZ16} F.~Fodor, V.~V\'igh and T.~Zarn\'ocz, \emph{On the angle sum of lines}, Arch.\ Math.\
  \textbf{106} (2016), 91--100.
\end{thebibliography}

\end{document}
